\documentclass[10pt,a4paper]{amsart}
\usepackage[margin=19mm]{geometry}
\usepackage{amsmath,amssymb,mathtools}
\usepackage[scr=rsfs,cal=euler]{mathalfa}
\usepackage{xfrac}
\usepackage[unicode,hidelinks,bookmarks=false]{hyperref}
\newcommand{\ie}{i.e.\ }
\newcommand{\cf}{cf.\ }
\newcommand{\resp}{respectively\ }
\newcommand{\Subsection}{\subsection}
\providecommand{\nobreakdash}{\nobreak\hbox{-}}
\setlength{\emergencystretch}{2em}
\multlinegap0pt
\usepackage{bm}
\usepackage{bbm}
\usepackage{dsfont}
\usepackage{xspace}
\usepackage{cancel}
\usepackage{mathtools}
\usepackage{amsthm}
\makeatletter
\@ifundefined{theorem}{\newtheorem{theorem}{Theorem}}{}
\makeatother
\title[Controllability and mixing for acoustic wave motions]{Controllability and mixing for acoustic wave motions}
\author{Zhe Jiao, Xiao Li, Qin Zhao}
\date{Source: arXiv:2508.16896v1 (CC BY 4.0)}
\begin{document}
\maketitle

\noindent\emph{Source and licence.} Controllability and mixing for acoustic wave motions. Version arXiv:2508.16896v1 (CC BY 4.0), distributed under the Creative Commons Attribution 4.0 International licence, as recorded in the arXiv licence metadata for this exact version and shown on the abstract page https://arxiv.org/abs/2508.16896v1. Full rights notice: licenses/arxiv-2508.16896-CC-BY-4.0.txt.\par\smallskip

\noindent\textit{Editorial scope.} Counted unit: the Theorem whose source label is \texttt{controllability} in the article (source0.tex lines 1087--1093, source label \texttt{controllability}), delivered together with the article's own complete original proof of it (source0.tex lines 1094--1160, one \texttt{proof} environment of 67 lines). No mathematics is written, repaired or re-derived: statement and proof are the article's own text line for line. Required context retained uncounted: 1 (lines 171--173); 2 (lines 175--183); 4 (lines 190--192); adjoint (lines 600--610); euler-lagrange (lines 1077--1083). Every definition, display and label of those lines is retained unchanged. Omitted from this package and not redistributed: 1--170, 174--174, 184--189, 193--599, 611--1076, 1084--1086, 1161--1297. Nothing inside a retained excerpt is omitted or altered: every retained body is delivered exactly as the article's lines are written, so no omission receipt of any kind accompanies this package. Citations: the retained excerpts carry no citation command at all, so no bibliography entry of the article is transcribed and no reference is redistributed. Reference adaptation: none was needed and none was made; every reference inside the delivered unit resolves inside it, so no unresolved reference or citation is produced. Printed slips kept literal: no printed word, symbol, hypothesis or number of the counted statement or of its proof was altered. Layout and class: the article's own class is \texttt{elsarticle} with packages including amssymb, amsfonts, verbatim, xcolor, amsmath, amsthm, extarrows, geometry, graphicx, booktabs, adjustbox, multirow, lineno, hyperref; the wrapper is the lane's pinned \texttt{amsart} editorial wrapper at 10pt on A4, which restates the theorem-like environments the retained text needs and adds the article's own macro definitions that the retained text needs (none), with extra wrapper packages (bm, bbm, dsfont, xspace, cancel, mathtools). The delivered numbering is the unit's own. Assets: no retained span contains a figure environment, a graphics-inclusion command, a PDF-page inclusion, a table or a TikZ picture; no image file is copied into this package, the package declares no asset dependency and no image file is redistributed with it. Licence: version arXiv:2508.16896v1 of this article is distributed under the Creative Commons Attribution 4.0 International licence, as recorded in the arXiv licence metadata for this exact version and shown on the abstract page https://arxiv.org/abs/2508.16896v1. A full rights notice is delivered at licenses/arxiv-2508.16896-CC-BY-4.0.txt. Characters: every retained excerpt is pure ASCII, so the delivered engine, XeLaTeX with its default Latin Modern text font, reports no missing character.
\begin{equation}
	 u_{tt} = c^2\triangle u, \quad \textrm{in $\mathbb{R}^{+} \times \Omega$.} \label{1}
\end{equation}

\begin{equation} \label{2}
\left
    \{
        \begin{array}{ll}
            -\rho u_t= m v_{tt} - \mathrm{div}_{\Gamma}\left(\sigma\nabla_{\Gamma} v\right) + d v_t + k v + f, \quad \textrm{on $\mathbb{R}^{+} \times \Gamma_1$},\\ 
            \partial_{\mathbf{n}}u =v_t, \quad \textrm{on $\mathbb{R}^{+} \times \Gamma_1$}, \quad \partial_{\mathbf{n}}u = 0, \quad \textrm{on $\mathbb{R}^{+} \times \Gamma_0$}.
        \end{array} 
\right. 
\end{equation}

\begin{equation}\label{4}
u(0, x)=u_0(x),\quad u_t(0, x)=u_1(x),\quad v(0, x)=v_0(x),\quad v_t(0, x)=v_1(x). 
\end{equation}

\begin{equation} \label{adjoint}
\left
    \{
        \begin{array}{ll}
        		\phi_{tt} = c^2\triangle\phi, \quad \textrm{in $\mathbb{R}^{+} \times \Omega$},\\
            -\rho \phi_t= m \delta_{tt} - \mathrm{div}_{\Gamma}\left(\sigma\nabla_{\Gamma} \delta\right) - d \delta_t + k \delta, \quad \textrm{on $\mathbb{R}^{+} \times \Gamma_1$},\\ 
            \partial_{\mathbf{n}}\phi =\delta_t, \quad \textrm{on $\mathbb{R}^{+} \times \Gamma_1$}, \quad \partial_{\mathbf{n}}\phi = 0, \quad \textrm{on $\mathbb{R}^{+} \times \Gamma_0$},\\
            \left(\phi(T, \cdot), \delta(T, \cdot), \phi_t(T, \cdot), \delta_t(T, \cdot)\right) = (\phi_0, \delta_0, \phi_1, \delta_1), \quad \textrm{in $ \Omega$}.
        \end{array} 
\right.
\end{equation}

\begin{equation}\label{euler-lagrange}
\begin{split}
&\int_{0}^{T}\int_{\Gamma_1}\left(\hat{\delta}_t\delta_t+ \zeta^2\hat{\delta}_{tt}\delta_{tt} \right) \mathrm{d}\Gamma \mathrm{d}t + \int_{\Omega}\rho u_{t}(0)\phi_t(0)\mathrm{d}x + \int_{\Gamma_1}m c^2 v_{t}(0)\delta_t(0)\mathrm{d}\Gamma\\
&\quad - \int_{\Omega}\rho u(0)\phi_{tt}(0)\mathrm{d}\Gamma - \int_{\Gamma_1}\rho c^2 v(0)\phi_t(0)\mathrm{d}\Gamma + \int_{\Gamma_1}\rho c^2 u(0)\delta_t(0)\mathrm{d}\Gamma\\
&\quad + \int_{\Gamma_1}c^2 d v(0)\delta_t(0)\mathrm{d}\Gamma - \int_{\Gamma_1}m c^2 v(0)\delta_{tt}(0)\mathrm{d}\Gamma= 0
\end{split}
\end{equation}

\begin{theorem}\label{controllability}
Suppose that $(\hat{\phi}, \hat{\delta}, \hat{\phi}_t, \hat{\delta}_t)$ is the minimizer of the functional $\mathcal{J}(\cdot)$. Then the control function 
\[
f = \frac{1}{c^2}\left[\hat{\delta_t} - (\zeta^2\hat{\delta}_{tt})_t\right]
\]
ensures the exact controllability of problem~\eqref{1}-\eqref{4}.
\end{theorem}

\begin{proof}
Let $(u, v)$ and $(\phi, \delta)$ be the strong solutions of problem~\eqref{1}-\eqref{4} and~\eqref{adjoint}, respectively. Assume that $(\phi, \delta)$ is smooth enough. Then by integrating by parts, we deduce from the first equation of~\eqref{adjoint} that
\begin{equation*}
\begin{split}
0=&\int_{0}^{T}\int_{\Omega}\rho u\left( {\phi}_{ttt} - c^2\triangle{\phi_t}\right)\mathrm{d}x \mathrm{d}t\\
=&\left[\int_{\Omega}\rho u {\phi}_{tt} \mathrm{d}x\right]_{t=0}^{t=T} - \left[\int_{\Omega}\rho u_{t} {\phi_t}\mathrm{d}x \right]_{t=0}^{t=T} + \int_{0}^{T}\int_{\Omega}\rho c^2\left( \triangle u{\phi_t} - u\triangle{\phi_t}\right)\mathrm{d}x \mathrm{d}t\\
=&\left[\int_{\Omega}\rho u {\phi}_{tt} \mathrm{d}x\right]_{t=0}^{t=T} - \left[\int_{\Omega}\rho u_{t} {\phi_t}\mathrm{d}x \right]_{t=0}^{t=T} + \int_{0}^{T}\int_{\Gamma}\rho c^2 \left(\partial_{\mathbf{n}}u{\phi_t} - u\partial_{\mathbf{n}}{\phi_t} \right) \mathrm{d}\Gamma \mathrm{d}t,
\end{split}
\end{equation*}
which implies from the boundary conditions that
\begin{equation} \label{explicit_1}
\left[\int_{\Omega}\rho u {\phi}_{tt} \mathrm{d}x\right]_{t=0}^{t=T} - \left[\int_{\Omega}\rho u_{t} {\phi_t}\mathrm{d}x \right]_{t=0}^{t=T} + \int_{0}^{T}\int_{\Gamma_1}\rho c^2 \left(v_{t}{\phi_t} - u{\delta}_{tt} \right) \mathrm{d}\Gamma \mathrm{d}t=0.
\end{equation}
Similarly, we also deduce from the second equation of~\eqref{adjoint} and~\eqref{2} that
\begin{equation}\label{explicit_2}
\begin{split}
0=&\int_{0}^{T}\int_{\Gamma_1}c^{2}v\left( \rho {\phi}_{tt} + m {\delta}_{ttt} - \mathrm{div}_{\Gamma}\left(\sigma\nabla_{\Gamma} {\delta_t}\right) - d {\delta}_{tt} + k {\delta_t} \right)\mathrm{d}\Gamma \mathrm{d}t\\
=&\left[\int_{\Gamma_1}\rho c^{2}v  {\phi_t}\mathrm{d}\Gamma\right]_{t=0}^{t=T} - \int_{0}^{T}\int_{\Gamma_1}\rho c^{2}v_{t}  {\phi_t}\mathrm{d}\Gamma \mathrm{d}t + \int_{0}^{T}\int_{\Gamma_1}m c^{2}v  {\delta}_{ttt}\mathrm{d}\Gamma \mathrm{d}t \\
&+\int_{0}^{T}\int_{\Gamma_1}c^{2}v\left[- \mathrm{div}_{\Gamma}\left(\sigma\nabla_{\Gamma} {\delta_t}\right)  - d {\delta}_{tt} + k {\delta_t} \right]\mathrm{d}\Gamma \mathrm{d}t
\end{split}
\end{equation}
and
\begin{equation} \label{explicit_3}
\begin{split}
 &\int_{0}^{T}\int_{\Gamma_1}m c^{2}v {\delta}_{ttt}\mathrm{d}\Gamma \mathrm{d}t \\
 &= \left[\int_{\Gamma_1}m c^{2}v  {\delta}_{tt}\mathrm{d}\Gamma\right]_{t=0}^{t=T} - \left[\int_{\Gamma_1}m c^{2}v_{t}  {\delta_t}\mathrm{d}\Gamma\right]_{t=0}^{t=T} + \int_{0}^{T}\int_{\Gamma_1}m c^{2}v_{tt}  {\delta_t}\mathrm{d}\Gamma \mathrm{d}t\\
 &= \left[\int_{\Gamma_1}m c^{2}v  {\delta}_{tt}\mathrm{d}\Gamma\right]_{t=0}^{t=T} - \left[\int_{\Gamma_1}m c^{2}v_{t}  {\delta_t}\mathrm{d}\Gamma\right]_{t=0}^{t=T} \\
 &\quad+ \int_{0}^{T}\int_{\Gamma_1} c^{2}  {\delta_t}\left[ -\rho u_{t} + \mathrm{div}_{\Gamma}\left(\sigma\nabla_{\Gamma} {v}\right) -d v_{t} - k v - f\right]\mathrm{d}\Gamma \mathrm{d}t\\
  &=  \left[\int_{\Gamma_1}m c^{2}v  {\delta}_{tt}\mathrm{d}\Gamma\right]_{t=0}^{t=T} - \left[\int_{\Gamma_1}m c^{2}v_{t}  {\delta_t}\mathrm{d}\Gamma\right]_{t=0}^{t=T} \\
  &\quad -\left[\int_{\Gamma_1}\rho c^{2}  {\delta_t} u\mathrm{d}\Gamma\right]_{t=0}^{t=T} + \int_{0}^{T}\int_{\Gamma_1}\rho c^{2}  {\delta}_{tt} u\mathrm{d}\Gamma \mathrm{d}t - \int_{0}^{T}\int_{\Gamma_1}c^{2}\delta_t f\mathrm{d}\Gamma \mathrm{d}t\\
 &\quad -\left[\int_{\Gamma_1} c^{2} d {v}\delta_t\mathrm{d}\Gamma \right]_{t=0}^{t=T}+ \int_{0}^{T}\int_{\Gamma_1} c^{2}  {v}\left[ \mathrm{div}_{\Gamma}\left(\sigma\nabla_{\Gamma} {\delta_t}\right) + d \delta_{tt} - k \delta_t\right]\mathrm{d}\Gamma \mathrm{d}t.
 \end{split}
\end{equation}

Combing~\eqref{explicit_1}, \eqref{explicit_2} and~\eqref{explicit_3} yields
\begin{equation}\label{explicit_4}
\begin{split}
&\left[\int_{\Omega}\rho u {\phi}_{tt} \mathrm{d}x\right]_{t=0}^{t=T} - \left[\int_{\Omega}\rho u_{t} {\phi_t}\mathrm{d}x \right]_{t=0}^{t=T} + \left[\int_{\Gamma_1}\rho c^{2}v  {\phi_t}\mathrm{d}\Gamma\right]_{t=0}^{t=T}\\
&+\left[\int_{\Gamma_1}m c^{2}v  {\delta}_{tt}\mathrm{d}\Gamma\right]_{t=0}^{t=T} - \left[\int_{\Gamma_1}m c^{2}v_{t}  {\delta_t}\mathrm{d}\Gamma\right]_{t=0}^{t=T}-\left[\int_{\Gamma_1}\rho c^{2}  {\delta_t} u\mathrm{d}\Gamma\right]_{t=0}^{t=T}\\
&-\left[\int_{\Gamma_1} c^{2} d {v}\delta_t\mathrm{d}\Gamma \right]_{t=0}^{t=T} - \int_{0}^{T}\int_{\Gamma_1}c^{2}\delta_t f\mathrm{d}\Gamma \mathrm{d}t=0.
 \end{split}
\end{equation}
Plugging $f = -\frac{1}{c^2}\left(\hat{\delta_t} - (\zeta^2\hat{\delta}_{tt})_t\right)$ into~\eqref{explicit_4} and integrating by parts, we obtain
\begin{equation}\label{explicit_5}
\begin{split}
&\left[\int_{\Omega}\rho u {\phi}_{tt} \mathrm{d}x\right]_{t=0}^{t=T} - \left[\int_{\Omega}\rho u_{t} {\phi_t}\mathrm{d}x \right]_{t=0}^{t=T} + \left[\int_{\Gamma_1}\rho c^{2}v  {\phi_t}\mathrm{d}\Gamma\right]_{t=0}^{t=T}\\
&+\left[\int_{\Gamma_1}m c^{2}v  {\delta}_{tt}\mathrm{d}\Gamma\right]_{t=0}^{t=T} - \left[\int_{\Gamma_1}m c^{2}v_{t}  {\delta_t}\mathrm{d}\Gamma\right]_{t=0}^{t=T}-\left[\int_{\Gamma_1}\rho c^{2}  {\delta_t} u\mathrm{d}\Gamma\right]_{t=0}^{t=T}\\
&-\left[\int_{\Gamma_1} c^{2} d {v}\delta_t\mathrm{d}\Gamma \right]_{t=0}^{t=T}+ \int_{0}^{T}\int_{\Gamma_1} \left(\hat{\delta}_{t}\delta_t  + \zeta^2\hat{\delta}_{tt} \delta_{tt}\right)\mathrm{d}\Gamma \mathrm{d}t=0.
 \end{split}
\end{equation}
Subtracting~\eqref{explicit_5} from~\eqref{euler-lagrange} gives
\begin{equation}\label{explicit_6}
\begin{split}
&\int_{\Omega}\rho u_{t}(T)\phi_t(T)\mathrm{d}x + \int_{\Gamma_1}m c^2 v_{t}(T)\delta_t(T)\mathrm{d}\Gamma- \int_{\Omega}\rho u(T)\phi_{tt}(T)\mathrm{d}x - \int_{\Gamma_1}\rho c^2 v(T)\phi_t(T)\mathrm{d}\Gamma\\
&  + \int_{\Gamma_1}\rho c^2 u(T)\delta_t(T)\mathrm{d}\Gamma+ \int_{\Gamma_1}c^2 d v(T)\delta_t(T)\mathrm{d}\Gamma - \int_{\Gamma_1}m c^2 v(T)\delta_{tt}(T)\mathrm{d}\Gamma = 0.
\end{split}
\end{equation}
Choosing $\phi_1=0$ and $\delta_1=0$, one easily deduces 
\[
\int_{\Omega}\rho u(T)\phi_{tt}(T)\mathrm{d}x + \int_{\Gamma_1}m c^2 v(T)\delta_{tt}(T)\mathrm{d}\Gamma = 0, \quad \forall (\phi_0, \delta_0),
\]
which implies $u(T) = 0$ and $v(T) = 0$. Furthermore, the equation \eqref{explicit_6} reduces to
\[
\int_{\Omega}\rho u_{t}(T)\phi_t(T)\mathrm{d}x + \int_{\Gamma_1}m c^2 v_{t}(T)\delta_t(T)\mathrm{d}\Gamma = 0, \quad \forall (\phi_0, \delta_0).
\]
Hence, $(u_{t}(T), v_{t}(T))=(0, 0)$. This concludes the proof of Theorem~\ref{controllability}.
\end{proof}

\end{document}
